\section{AI-Mediated Transformations and Mobile Lineage}
\label{sec:ai}

AI-assisted development is not one transformation. It ranges from completion of a few lines to generation of an application skeleton, selection of dependencies, modification across multiple files, test execution, debugging, and release preparation. The relevant provenance question is therefore not merely whether ``AI was used.'' It is which artifact-changing decisions were delegated, which context influenced them, and how the resulting outputs were bound to the final application.

\subsection{From text descriptions to mobile artifacts}

Text2App's July 2021 v2 translates descriptions through a Simplified App Representation and predefined MIT App Inventor templates to an APK, not unrestricted native Android synthesis~\cite{text2app}. Its synthetic corpus and volunteer evaluation use different outcomes; exact-SAR and functional judgments have unreconciled 111/112 denominators. Compilation does not certify security, and templates and literal processing belong in the accepted-output history.

AppForge evaluates generated Android projects and agents, including compilation-feedback refinement, not authenticated lineage~\cite{appforge}. Functional success counts compile-plus-all-tests over all tasks; test-pass and crash rates condition on compiled apps, and no joint outcome is supplied to infer crash rates among functional successes. Neither success nor a final diff certifies security or records prompts, retrieval, tool state and accepted revisions.

These systems do not make traditional repackaging evidence obsolete. A generated app can itself be copied, repackaged, or combined with existing code. What changes is the space of legitimate relatedness. Multiple authorized outputs may descend from the same prompt, template, or model-mediated plan without one being copied from another. Direction cannot be inferred from pairwise similarity alone.

\subsection{Nondeterminism and reproducibility}

Conventional build reproducibility assumes that a declared input set and deterministic procedure should yield the same bytes. Generative systems may sample different outputs for the same prompt, call external tools, or depend on changing model and retrieval services. Exact reproduction can therefore be unavailable even when the process is authorized.

This does not eliminate reproducibility as an objective; it changes what must be reproduced. One target is \emph{record reproducibility}: an independent verifier can confirm the prompt or specification reference, model or service identity, parameters, tool sequence, retrieved materials, and accepted outputs. Another is \emph{decision replay}: a recorded output is retained and subsequent deterministic build steps reproduce the APK. A third is \emph{distributional replication}: repeated generations satisfy a defined property, which is an empirical model claim rather than artifact lineage.

For mobile supply-chain assurance, preserving the accepted generated source or patch is usually more important than reproducing the model's sampling event. The build attestation should treat that accepted output as a material with a digest. If policy requires regeneration, the record must acknowledge that regenerated code may be a sibling rather than an identical ancestor.

\subsection{Security of generated code}

Generation provenance and secure behavior are separate claims. Pearce et al. tested targeted security scenarios using the 2021 Copilot preview; Perry et al. studied 47 participants with codex-davinci-002; Sandoval et al. studied 58 students implementing a C API~\cite{pearce,perry,sandoval}. Their task designs, models and vulnerability oracles differ. Pearce's designated-CWE checks do not establish functional correctness; Perry's task outcomes are mixed; Sandoval's compiling- versus passing-function comparisons and chosen non-inferiority tolerance answer different questions. Khoury et al.'s 2026 revision still concerns historical ChatGPT-3.5 programs and selected attack classes~\cite{khoury}. These are not current Android security or universal repair rates.

The transferable design obligation is to analyze the accepted code and exact built artifact. An authorized generated change may be vulnerable, while an unauthorized derivative may remove a flaw. Provenance can identify a recorded change; property-specific analysis supplies the security evidence. Human acceptance likewise records an interaction rather than automatically establishing informed review. A review policy must specify the displayed change, checks, warnings and the approving principal's authority.

\subsection{Dependencies and package hallucination}

Generative systems can recommend dependencies by name. Spracklen et al. measure nonexistent package recommendations in Python and JavaScript generation, using install-command and follow-up-prompt heuristics against PyPI/npm snapshots~\cite{spracklen}. Their registry-abuse scenario motivates a dependency risk, but it does not measure an Android attack rate. Transfer to Maven coordinates, Gradle plugins, native packages or download URLs is a hypothesis requiring ecosystem-specific validation; JavaScript dependencies in hybrid apps are closer to the measured setting. The suggestion may enter source or build configuration even when no copied code appears in the model output.

The missing-name fraction counts recommendations, not programs, installs or victims~\cite{spracklen}; registry presence is not a trust oracle. The reported mitigation reduces one DeepSeek hallucination rate from 16.14\% to 2.66\% while HumanEval falls from 51.4 to 25.3. No package was published for the hypothesized abuse. Thus the measured name errors do not establish compromise or a cost-free mitigation.

This risk links AI assistance directly to supply-chain provenance. A robust record should capture the dependency request, resolver result, repository origin, version, integrity digest, and whether the dependency was already approved. A policy can reject unrecognized repositories or require lockfile review. Binary component analysis then checks whether the resolved material appears in the APK. Neither the prompt log nor the binary detector alone provides the full chain.

Package hallucination also motivates a counter-history for library attribution. A newly registered malicious package may intentionally implement an interface or name suggested by a model. The package's code is not necessarily similar to an established library, so clone detection may not flag it. Repository provenance, registration chronology, maintainer identity, and dependency-policy evidence become central.

\subsection{Training-data and privacy lineage}

CodexLeaks uses public-code search and author review to validate sensitive-output associations, not authenticated membership in an unavailable training corpus~\cite{codexleaks}. Its 43 confirmed leakage responses among 2,560 responses and the abstract's prompt-level percentage have different denominators. Prefilter limitations, short-output calibration and public-source adjudication bound the inference. Neither output similarity nor a prompt/model record identifies an exact training example or certifies present-day defenses.

A final APK therefore needs artifact-specific investigation before sensitive-output observations become a lineage or exposure claim. Context retention itself can expose secrets: retain access-controlled details and selective commitments that bind accepted output to the final artifact. Technical resemblance does not alone resolve licensing, confidentiality or permitted use.

\subsection{Multi-principal and tool-mediated authorship}

An agentic workflow can include a human requester, model provider, model instance, retrieval source, code repository, build service, package registry, test system, reviewer, and release signer. These are not equivalent principals. The model service may produce candidate text but lack release authority. A human may approve a patch but not control the signing key. A build service may emit an attestation but not own the source.

Supply-chain layouts are well suited to representing these roles when model-mediated steps are treated as explicit transformations. The layout can require that generated outputs pass tests, static checks, dependency policy, and human review before signing. It can also distinguish advisory output from accepted material. This is more precise than adding one boolean ``AI-generated'' field to the APK metadata.

The same role separation helps incident response. If a vulnerability is found, investigators can ask whether it originated in a generated patch, retrieved example, dependency recommendation, manual edit, or build-time transformation. The answer guides remediation and model or policy changes. Without step records, the final Git history may collapse several causes into one commit.

\begin{table}[t]
\caption{AI-mediated development risks and the evidence needed to reason about them.}
\label{tab:ai-risks}
\footnotesize
\begin{tabularx}{\textwidth}{L{2.35cm}Y Y L{1.6cm}}
\toprule
Risk or ambiguity & Why final-artifact similarity is insufficient & Complementary evidence & Main claim \\
\midrule
Sibling generations (analytical scenario)~\cite{text2app,appforge} & Similarity does not reveal whether one output copied another & Specification binding, generation records, accepted-output digests & \claimlevel{3}--\claimlevel{4} \\
Insecure generated code~\cite{pearce,perry,sandoval} & Origin says nothing about security of accepted behavior & Static/dynamic analysis, tests, review evidence, exact artifact binding & \claimlevel{6} \\
Hallucinated dependencies~\cite{spracklen} & Malicious package may be novel and not resemble a known library & Resolver provenance, repository identity, lockfile, component verification & \claimlevel{1}+\claimlevel{4} \\
Sensitive-output association~\cite{codexleaks} & Exact training membership and directed ancestry are not authenticated & Model/prompt record, public-source investigation, access-controlled context & Disclosure risk; \claimlevel{3} unresolved \\
Agentic multi-file edits~\cite{appforge} & A final diff alone need not preserve the intermediate history & Step attestations, tool transcripts, policy checkpoints & \claimlevel{4}--\claimlevel{5} \\
Nondeterministic regeneration & Exact output may not be reproducible from a prompt alone & Preserve accepted output; separate record replay from model replication & \claimlevel{0}+\claimlevel{4} \\
\bottomrule
\end{tabularx}
\end{table}

\subsection{A minimum evidence record for generated transformations}

A generated transformation record should identify: the parent source revision; task or specification identifier; model or service identity at the granularity available; generation parameters; tool and retrieval accesses; candidate-output digests; the accepted patch or file digests; tests and analyses run; human or policy approvals; and the downstream build product. Sensitive context can be access controlled, but the public or shared record should still bind the accepted output to the final APK.

This record serves three purposes. It establishes direction from a parent revision to an accepted generated change. It separates generation from approval and release. It gives artifact and behavior analyses exact objects to examine. The record does not prove that the model never used unrecorded training material or that the accepted code is secure. It makes those residual questions explicit rather than hiding them behind a generic provenance claim.
